\documentclass[final]{beamer}

%================================================
% POSTER PAGE SIZE
%================================================
%
% A0 landscape poster.
%
% scale=1.05 enlarges Beamer's default poster typography slightly.
%
% POSSIBLE CHANGES:
%   scale=1.00  -> slightly smaller overall
%   scale=1.05  -> current
%   scale=1.10  -> larger, with greater overflow risk
%
\usepackage[
  size=a0,
  orientation=landscape,
  scale=1.05
]{beamerposter}

\usetheme{default}

%================================================
% PACKAGES
%================================================

\usepackage{graphicx}
\usepackage{booktabs}
\usepackage{array}
\usepackage{xcolor}
\usepackage{amsmath}
\usepackage{ragged2e}

% Figures must be stored in the figures folder.
\graphicspath{{figures/}}

% Remove Beamer navigation symbols.
\setbeamertemplate{navigation symbols}{}

%================================================
% COLOURS
%================================================

\definecolor{projectblue}{RGB}{28,76,128}
\definecolor{lightblue}{RGB}{235,242,249}

\setbeamercolor{headline}{
  bg=projectblue,
  fg=white
}

\setbeamercolor{block title}{
  bg=projectblue,
  fg=white
}

\setbeamercolor{block body}{
  bg=lightblue,
  fg=black
}

%================================================
% MAIN FONT-SIZE CONTROLS
%================================================
%
% In \fontsize{A}{B}:
%
%   A = letter size
%   B = line spacing
%
% B should normally be larger than A.
%
%------------------------------------------------
% MAIN POSTER TITLE
%------------------------------------------------
%
% CURRENT:
%   98pt letters with 110pt line spacing.
%
% POSSIBLE CHANGES:
%
%   88/100   -> large and safe
%   94/106   -> very large
%   98/110   -> current
%   104/118  -> extremely large
%
\newcommand{\PosterTitleFont}{%
  \fontsize{98}{110}\selectfont
}

%------------------------------------------------
% AUTHOR / COLLEGE / SUPERVISOR LINE
%------------------------------------------------
%
% CURRENT:
%   50pt letters with 60pt line spacing.
%
% POSSIBLE CHANGES:
%
%   42/50 -> smaller
%   46/56 -> large
%   50/60 -> current
%   54/65 -> very large
%
\newcommand{\PosterSubtitleFont}{%
  \fontsize{50}{60}\selectfont
}

%------------------------------------------------
% CONTENT-BLOCK TITLES
%------------------------------------------------
%
% Controls:
%
%   Research Aim
%   Frozen Model Benchmark
%   Model Explainability
%   Reproducibility
%   Conclusions
%
% POSSIBLE CHANGES:
%
%   32/38 -> slightly smaller
%   34/40 -> current
%   38/45 -> larger
%   42/50 -> very large, increases vertical space
%
\setbeamerfont{block title}{
  series=\bfseries,
  size*={34pt}{40pt}
}

%------------------------------------------------
% NORMAL BODY TEXT
%------------------------------------------------
%
% POSSIBLE CHANGES:
%
%   \large -> smaller
%   \Large -> current
%   \LARGE -> larger, but substantial overflow risk
%
\setbeamerfont{block body}{
  size=\Large
}

\setbeamerfont{itemize/enumerate body}{
  size=\Large
}

\setbeamerfont{itemize/enumerate subbody}{
  size=\large
}

%================================================
% BULLET POSITIONING
%================================================
%
% These margins keep the triangular bullets inside the coloured blocks.
%
% POSSIBLE CHANGES:
%
%   1.8em -> tighter
%   2.1em -> current
%   2.4em -> further inward
%
\setlength{\leftmargini}{2.1em}
\setlength{\leftmarginii}{1.7em}

\setbeamertemplate{itemize item}{%
  \raisebox{0.15ex}{%
    \color{projectblue}
    \scriptsize$\blacktriangleright$%
  }%
}

%================================================
% BLOCK DESIGN
%================================================

\setbeamertemplate{block begin}{%

  % Small vertical space above each block.
  \vskip0.4ex

  %----------------------------------------------
  % DARK-BLUE BLOCK-TITLE BAR
  %----------------------------------------------
  %
  % sep controls bar thickness.
  %
  %   0.75ex -> thinner
  %   0.90ex -> current
  %   1.10ex -> thicker
  %
  \begin{beamercolorbox}[
    rounded=true,
    shadow=true,
    sep=0.9ex,
    leftskip=0.6ex,
    rightskip=0.6ex
  ]{block title}

    \usebeamerfont*{block title}
    \insertblocktitle

  \end{beamercolorbox}

  \ifbeamercolorempty[bg]{block body}
  {}
  {\nointerlineskip\vskip-0.4pt}

  \usebeamerfont{block body}

  %----------------------------------------------
  % PALE-BLUE BLOCK BODY
  %----------------------------------------------
  %
  % leftskip and rightskip create internal margins.
  %
  % Increase leftskip if bullets still look close to the edge.
  %
  \begin{beamercolorbox}[
    rounded=true,
    shadow=true,
    sep=0.9ex,
    leftskip=1.6ex,
    rightskip=1.2ex,
    vmode
  ]{block body}

  \raggedright
}

\setbeamertemplate{block end}{%
  \end{beamercolorbox}
  \vskip0.45ex
}

%================================================
% POSTER INFORMATION
%================================================

\title{
  Benchmarking ImageNet Transfer Learning for\\
  Restricted Dog Image Classification
}

\author{Ronan Downes}

\institute{
  MSc in Data Analytics, CCT College Dublin
  \qquad
  Supervisor: Sam Weiss
}

%================================================
% DOCUMENT
%================================================

\begin{document}

\begin{frame}[t]

%================================================
% MAIN TITLE BANNER
%================================================
%
% ht controls the height of the dark-blue header.
%
% POSSIBLE CHANGES:
%
%   0.150 -> shallower
%   0.160 -> current
%   0.170 -> deeper
%
% A deeper header gives the large title more room but reduces the
% vertical space available for the three content columns.
%
\begin{beamercolorbox}[
  wd=\paperwidth,
  ht=0.160\paperheight,
  dp=1.25cm,
  center
]{headline}

  {\PosterTitleFont
   \bfseries
   \inserttitle
   \par}

  % Space between main title and the author/institution line.
  \vspace{0.45cm}

  {\PosterSubtitleFont
   \bfseries
   \insertauthor
   \hspace{1.5em}
   \insertinstitute}

\end{beamercolorbox}

% Gap between the blue header and the poster content.
\vspace{0.30cm}

%================================================
% MAIN POSTER COLUMNS
%================================================
%
% IMPORTANT:
%
% The two narrow empty columns are deliberate white gutters.
% They stop the first three blue title bars from appearing to form
% one continuous bar.
%
% WIDTHS:
%
%   Left content column    = 0.305
%   First white gutter     = 0.0125
%   Centre content column  = 0.325
%   Second white gutter    = 0.0125
%   Right content column   = 0.305
%
% Total explicitly occupied width = 0.960 paper width.
%
% The whole column region is 0.98 paper width, leaving a little
% additional outside breathing room.
%
% GUTTER OPTIONS:
%
%   0.008  -> narrow separation
%   0.0125 -> current clear separation
%   0.016  -> wider separation
%
% If you increase the gutters, reduce the three content columns.
%
\begin{columns}[t,totalwidth=0.98\paperwidth]

%================================================
% LEFT COLUMN
%================================================

\begin{column}{0.305\paperwidth}

\begin{block}{Research Aim}

To benchmark six ImageNet-pretrained CNNs for
restricted-versus-unrestricted dog-image classification and evaluate
the effect of limited fine-tuning.

\end{block}

\begin{block}{Dataset and Experimental Design}

\begin{itemize}

  % Vertical distance between bullet items.
  %
  % 0.10em -> tighter
  % 0.18em -> current
  % 0.28em -> more open
  %
  \setlength{\itemsep}{0.18em}
  \setlength{\parskip}{0pt}

  \item Stanford Dogs image dataset: \textbf{30 breeds}.

  \item \textbf{6 restricted} and \textbf{24 unrestricted} breeds.

  \item Fixed training, validation and held-out test splits.

  \item Six ImageNet-pretrained CNNs benchmarked as frozen feature
  extractors.

  \item The selected model was fine-tuned by unfreezing its final
  20 layers.

  \item Evaluation used accuracy, precision, recall, F1, ROC-AUC,
  PR-AUC, MCC and confusion counts.

\end{itemize}

\vspace{0.2cm}

\centering

% Pipeline diagram size.
%
% POSSIBLE HEIGHTS:
%
%   0.15 -> smaller
%   0.17 -> current
%   0.19 -> larger
%
\includegraphics[
  width=0.98\linewidth,
  height=0.17\paperheight,
  keepaspectratio
]{training_pipeline_diagram.pdf}

\end{block}

\begin{block}{Six Benchmark Architectures}

\centering
\textbf{VGG16 $\cdot$ ResNet50 $\cdot$ InceptionV3}\\[0.12cm]
\textbf{Xception $\cdot$ InceptionResNetV2 $\cdot$ NASNetMobile}

\end{block}

\begin{block}{Reproducibility}

% The filenames use a slightly smaller font because they are long.
%
% Change \large to \Large if there is enough width.
%
\begin{itemize}
  \setlength{\itemsep}{0.12em}
  \item Data preparation
  \item Frozen-model benchmark
  \item Selected-model fine-tuning
  \item Grad-CAM explanation
\end{itemize}

The archive retains the four notebooks, figures, result tables,
predictions, checkpoints and complete \LaTeX{} source.

\end{block}

\begin{block}{Ethical Boundary}

\centering

% Important statement size.
%
%   25/31 -> smaller
%   27/33 -> current
%   30/36 -> larger
%
\fontsize{27}{33}\selectfont
\bfseries

High predictive performance does not justify autonomous enforcement.

\end{block}

\end{column}

%================================================
% WHITE GUTTER BETWEEN LEFT AND CENTRE
%================================================
%
% This empty column is intentional.
%
% Increase 0.0125 to 0.016 for a wider visual gap.
% Do not place content inside it.
%
\begin{column}{0.0125\paperwidth}
\end{column}

%================================================
% CENTRE COLUMN
%================================================

\begin{column}{0.325\paperwidth}

\begin{block}{Frozen Model Benchmark}

\centering

% Frozen-model comparison chart.
%
%   0.19  -> smaller
%   0.205 -> current
%   0.22  -> larger
%
\includegraphics[
  width=\linewidth,
  height=0.205\paperheight,
  keepaspectratio
]{model_zoo_four_metric_comparison_presentation.pdf}

\vspace{0.22cm}

\raggedright

\textbf{Best frozen model: InceptionResNetV2}\\[0.1cm]
97.08\% accuracy $\cdot$ 98.42\% recall $\cdot$ 92.57\% F1

\end{block}

\begin{block}{Fine-Tuning Result}

\centering

% Table-row spacing.
%
%   1.05 -> tighter
%   1.17 -> current
%   1.30 -> more open
%
\renewcommand{\arraystretch}{1.17}

\begin{tabular}{lcc}

\toprule

\textbf{Metric}
&
\textbf{Frozen}
&
\textbf{Fine-tuned}
\\

\midrule

Accuracy
&
97.08\%
&
\textbf{98.25\%}
\\

Precision
&
87.38\%
&
\textbf{93.88\%}
\\

Recall
&
\textbf{98.42\%}
&
96.84\%
\\

F1
&
92.57\%
&
\textbf{95.34\%}
\\

MCC
&
0.9101
&
\textbf{0.9428}
\\

\bottomrule

\end{tabular}

\vspace{0.28cm}

\raggedright

\textbf{Key trade-off:}
\textbf{15 fewer false positives}, but \textbf{3 additional false negatives}.

Fine-tuning improved overall balance and precision, but slightly
reduced recall.

\end{block}

\begin{block}{Fine-Tuned Confusion Matrix}

\centering

% Confusion-matrix size.
%
% WIDTH:
%   0.78 -> smaller
%   0.84 -> current
%   0.90 -> larger
%
% HEIGHT:
%   0.15 -> smaller
%   0.17 -> current
%   0.19 -> larger
%
\includegraphics[
  width=0.84\linewidth,
  height=0.17\paperheight,
  keepaspectratio
]{FineTuned_InceptionResNetV2_confusion_matrix.png}

\end{block}

\end{column}

%================================================
% WHITE GUTTER BETWEEN CENTRE AND RIGHT
%================================================

\begin{column}{0.0125\paperwidth}
\end{column}

%================================================
% RIGHT COLUMN
%================================================

\begin{column}{0.305\paperwidth}

\begin{block}{Model Explainability}

\centering

% German Shepherd Grad-CAM image.
%
%   0.105 -> smaller
%   0.115 -> current
%   0.125 -> larger
%
\includegraphics[
  width=0.95\linewidth,
  height=0.115\paperheight,
  keepaspectratio
]{TP_german_shepherd_n02106662_13368_gradcam.png}

\vspace{0.06cm}

{\Large\bfseries True positive: German Shepherd}

\vspace{0.12cm}

\raggedright

Attention concentrated around the head and ears, which contributed
strongly to the restricted prediction.

\vspace{0.22cm}

\centering

% Weimaraner false-positive Grad-CAM image.
\includegraphics[
  width=0.95\linewidth,
  height=0.115\paperheight,
  keepaspectratio
]{FP_weimaraner_n02092339_3028_gradcam.png}

\vspace{0.06cm}

{\Large\bfseries False positive: Weimaraner}

\vspace{0.12cm}

\raggedright

Body shape, pose, coat contrast and contextual structure may have
contributed to the incorrect restricted prediction.

\vspace{0.18cm}

\textbf{Interpretive boundary:}
Grad-CAM heatmaps are post-hoc visual diagnostics. They indicate
regions associated with a prediction, but they do not provide causal,
behavioural or legal explanations.

\end{block}

\begin{block}{Conclusions}

\begin{itemize}
  \setlength{\itemsep}{0.18em}
  \setlength{\parskip}{0pt}

  \item Six frozen ImageNet CNNs were benchmarked under one pipeline.

  \item InceptionResNetV2 produced the strongest frozen result.

  \item Fine-tuning improved precision and overall balance but reduced recall.

  \item The classifier cannot establish legal identity, ancestry,
  behaviour or dangerousness.

\end{itemize}

\end{block}

\end{column}

\end{columns}

\end{frame}

\end{document}
